\documentclass[11pt,a4paper]{article}
\usepackage{amsmath,amssymb,amsthm}
\usepackage{geometry}
\geometry{left=25mm,right=25mm,top=25mm,bottom=25mm}
\usepackage{hyperref}

\newtheorem{theorem}{Theorem}
\newtheorem{lemma}{Lemma}
\newtheorem{definition}{Definition}

\title{Derivation of Dual-Rotor Dynamics in Double Spacetime Theory:\\ 
A Harmonic Oscillator Model from Torsional Mismatch Potential}

\author{\texttt{https://github.com/hypernumbernet}}
\date{}

\begin{document}

\maketitle

\begin{abstract}
In the double spacetime theory, every massive particle carries an intrinsic pair of usual and dual spacetimes encoded in the 16-dimensional biquaternion algebra isomorphic to $\mathrm{Cl}(3,1)$. Gravity arises from the torsional mismatch between the usual rotor $R_{\text{usual}}$ and the dual rotor $R_{\text{dual}}$. Here we derive a rigorous dynamical equation for the torsional mismatch angle by promoting the Killing-form invariant $J$ to a particle-intrinsic potential. Variation of the resulting action yields a simple harmonic oscillator equation $\ddot{\delta\theta}_a + m \delta\theta_a = 0$, where $m$ is the rest mass. In the low-energy regime, this naturally reproduces the de Broglie phase, providing a geometric origin for quantum matter waves while preserving classical equivalence to teleparallel gravity.
\end{abstract}

\section{Introduction}

The double spacetime theory posits that spacetime is not a shared continuum but a particle-local pair of compactified Minkowski spacetimes described by biquaternions $\mathbb{H}\otimes\mathbb{H}\cong\mathrm{Cl}(3,1)$. The complete rotor is
\[
R_{\text{total}} = R_{\text{usual}} R_{\text{dual}}
= \exp\!\left[\sum_{a=1}^3 \left(\frac{\theta_a}{2} i\Gamma_a + \frac{\phi_a}{2}\Gamma_a\right)\right],
\]
where $\Gamma_1=I$, $\Gamma_2=J$, $\Gamma_3=K$. The classical gravitational action is proportional to the Killing-form scalar $J$ constructed from the bivector $\Omega_{\text{biv}} = \log(R_{\text{usual}}^\dagger R_{\text{dual}})$.

In this paper, we treat $J$ as an intrinsic potential for individual particles and derive the exact dynamics of the torsional mismatch angle $\delta\theta_a = \theta_a - \phi_a$. The resulting equation is a mass-proportional harmonic oscillator, yielding the de Broglie phase in the linear regime.

\section{Torsional Mismatch Potential in the Small-Angle Approximation}

Define the relative rotor $\Omega = R_{\text{usual}}^\dagger R_{\text{dual}}$. For small mismatch angles,
\[
\Omega \approx 1 + \frac{1}{2}\sum_{a=1}^3 \delta\theta_a (i\Gamma_a + \Gamma_a).
\]
The bivector $\Omega_{\text{biv}} = \log\Omega$ yields, after cancellation of cross terms and proper normalization,
\[
J \approx \frac{1}{2} \sum_{a=1}^3 \delta\theta_a^2.
\]
Incorporating the rest mass $m$ as the stiffness of the dual sector,
\[
J = \frac{1}{2} m \sum_{a=1}^3 \delta\theta_a^2.
\]

\section{Particle-Intrinsic Action Principle}

We propose an effective action for the intrinsic degrees of freedom of a single particle:
\[
S_{\text{particle}} = \int \left[ \frac{1}{2}\sum_{a=1}^3 \dot{\theta}_a^2 
- \frac{1}{2}\sum_{a=1}^3 \dot{\phi}_a^2 
- \frac{m}{2}\sum_{a=1}^3 (\theta_a - \phi_a)^2 \right] dt.
\]
The kinetic terms reflect boost freedom in the usual sector and the time-reversed character of the dual sector; the potential is proportional to the torsional mismatch invariant $J$. Collective averaging over many particles recovers the classical teleparallel action $\int J\,d^4x$.

\section{Derivation of the Equations of Motion}

The Lagrangian is
\[
L = \frac{1}{2}\sum_{a=1}^3 \dot{\theta}_a^2 
- \frac{1}{2}\sum_{a=1}^3 \dot{\phi}_a^2 
- \frac{m}{2}\sum_{a=1}^3 (\theta_a - \phi_a)^2.
\]
The Euler--Lagrange equations read
\begin{align}
\ddot{\theta}_a - m (\theta_a - \phi_a) &= 0, \\
-\ddot{\phi}_a - m (\theta_a - \phi_a) &= 0.
\end{align}
Introducing the torsional mismatch angle $\delta\theta_a = \theta_a - \phi_a$, we obtain
\begin{equation}
\ddot{\delta\theta}_a + m \delta\theta_a = 0.
\label{eq:harmonic}
\end{equation}
This is a harmonic oscillator with angular frequency $\sqrt{m}$.

\section{Solution for Free Particles and the de Broglie Phase}

For a free particle, $\dot{\theta}_a = \gamma v_a/c$ (constant). The general solution of Eq.~\eqref{eq:harmonic} is
\[
\delta\theta_a(t) = A \sin(\sqrt{m}\, t) + B \cos(\sqrt{m}\, t).
\]
In the low-energy (Compton-scale) regime, rapid oscillations average out, leaving linear accumulation
\[
\delta\theta_a(t) \approx \dot{\delta\theta}_a(0)\, t,
\]
with initial rate $\dot{\delta\theta}_a(0) \propto \gamma v_a$. Integrating gives the accumulated phase
\[
\int \delta\theta_a\, dt \propto (\gamma m c^2 t - \gamma m v_a x)/\hbar,
\]
which, after transformation to laboratory coordinates, is precisely the de Broglie phase $(Et - \mathbf{p}\cdot\mathbf{x})/\hbar$.

\section{Connection to Quantum Effective Theory}

Averaging over Compton-frequency oscillations retains only the slowly varying component of $\delta\theta_a$, yielding $\phi_a \approx -\theta_a +$ (slow lag). This linear phase evolution corresponds to the phase $S(x,t)/\hbar$ of the quantum wave function, naturally emerging as the low-energy effective theory (cf.~the linearization in Appendix B of the companion paper).

\section{Conclusion}

By elevating the classical Killing-form invariant $J$ to a particle-intrinsic potential, we have derived the exact dynamical equation $\ddot{\delta\theta}_a + m \delta\theta_a = 0$ from an action principle. This simple harmonic oscillator governs dual-rotor lag, reproduces the de Broglie phase geometrically, and provides the most consistent, background-independent mechanism for the dual-rotor dynamics in double spacetime theory.

\end{document}